\section{Construction details}\label{app:construction}
\sectionstate{partial}{\texttt{docs/detector\_spec/CONSTRUCTION\_DOSSIER.md}; Stephan's notes (2026-11-24), the 3498-A
stack-up, the production Gerbers, bench run configurations, February and May 2026 logbooks. New 2026-11-25.}
\intent{Everything about how the chambers were built that a reader trying to build or understand one would want, and
that the body of the paper does not need.}

\subsection{The pad pattern}\label{app:pattern}
A two-dimensional pad readout joins each pad to an X strip or a Y strip through its own via. In the CAST
Micromegas~\cite{Abbon2007CAST} the X and Y strips run along the diagonals of the pad grid. In that layout each pad has
two neighbors on its own strip, and pads of the same strip that touch at a corner can be joined on the pad layer, which
halves the number of vias. However, the strips then end along the edges of a square at every length, and routing them
to the connectors is difficult.

In our pattern the X strips run along $x$ and the Y strips along $y$, and the pads alternate between the views in a
checkerboard (\crefp{fig:chamber}{c}). A pad of strip X$_n$ then shares its four edges with pads of Y strips, and its
four corners with pads of the neighboring strips X$_{n-1}$ and X$_{n+1}$. No neighbor of a pad belongs to the same
strip. For this reason the charge of an avalanche is shared among as many strips as possible, and the 512 X and 512 Y
strips run straight to two edges of the board. The price is one via per pad, which is 262\,144 vias on a board.

\subsection{The readout board}\label{app:boardstack}
\Cref{tab:boardstack} gives the layers of the readout board as specified to the CERN MPT workshop (offer 3498-A). The
board has six copper layers (L3 to L8 in the Gerber files, where L1 and L2 are the artwork of the pillars and of the
resistive print). The layers add up to 528\,\textmu m, and the offer gives 530\,\textmu m. Under the active area,
L3 has no copper, and L7 carries only the land at the foot of each via.

\begin{table}[tbp]
  \centering\small
  \caption{Stack-up of the readout board, from the mesh side down (offer 3498-A of the CERN MPT workshop). Where a
  copper layer is given as two numbers, the second is the plating. $^\ast$: connector pads and the contacts of the resistive strips; no copper inside the active area.}
  \label{tab:boardstack}
  \begin{tabular}{lll}
    \hline
    layer & material & thickness \\
    \hline
    L3, top copper (contacts only)$^\ast$ & Cu, Ni/Au finish & 9 + 17\,\textmu m \\
    dielectric & UBE polyimide + EMC prepreg & 25 + 25\,\textmu m \\
    L4, pads & Cu & 17 + 17\,\textmu m \\
    dielectric & prepreg & 100\,\textmu m \\
    L5, Y strips & Cu & 17\,\textmu m \\
    dielectric & EMC prepreg & 100\,\textmu m \\
    L6, X strips & Cu & 17\,\textmu m \\
    dielectric & prepreg & 100\,\textmu m \\
    L7, via lands & Cu & 17 + 17\,\textmu m \\
    dielectric & EMC prepreg + UBE polyimide & 25 + 25\,\textmu m \\
    L8, bottom & no copper & \\
    \hline
    total & & 528\,\textmu m \\
    \hline
  \end{tabular}
\end{table}
\askStephan{Two disagreements with what we used so far: the insulator between the resistive strips and the pads is
50\,\textmu m here (our charge-spreading model assumed 50\,\textmu m Kapton + 25\,\textmu m glue), and the board is
0.53\,mm here (1.70\,mm in the mechanical CAD). Questions sent to Bertrand on 2026-11-24.}

\Cref{tab:material} gives the numbers drawn in \cref{fig:stack} for every layer of the chamber.
% Numbers: figs/out/material_stack.json (figs/make_material_stack.py, which also writes this tabular).
% Highland, per projection, electron momentum. Patterned layers averaged over their pattern.
\begin{table}[tbp]
  \centering\footnotesize
  \caption{The layers of a chamber at the center of the active area, in the order a lepton from the target crosses them,
  and the multiple scattering each layer causes alone: rms angle $\theta_0$ per projection for electrons of 4, 8 and
  16\,MeV kinetic energy (Highland formula). For a patterned layer the fill is the fraction of the area that it covers,
  and $x/X_0$ is averaged over the pattern. The summed rows apply the Highland formula to the summed $x/X_0$. The
  polyimide of the board is counted as glass-epoxy, and the aluminization of the window is ignored. $^\dagger$: thickness
  assumed. These are the numbers drawn in \cref{fig:stack}.}
  \label{tab:material}
  \setlength{\tabcolsep}{3.5pt}
  \begin{tabular}{llrrrrrr}
  \hline
   & & & & & \multicolumn{3}{c}{$\theta_0$ (\textdegree), $e^\pm$ of} \\
  layer & material & thickness & fill & $x/X_0$ (\%) & 4\,MeV & 8\,MeV & 16\,MeV \\
  \hline
  air, target to window & air & 200\,mm &  & 0.0658 & 3.24 & 1.70 & 0.87 \\
  entrance window, Mylar & Mylar & 60\,\textmu m &  & 0.0210 & 1.72 & 0.90 & 0.46 \\
  gas gap & Ar/iC$_4$H$_{10}$ 95/5 & 5\,mm &  & 0.0042 & 0.70 & 0.37 & 0.19 \\
  drift cathode, Kapton & Kapton & 50\,\textmu m &  & 0.0175 & 1.55 & 0.82 & 0.42 \\
  drift cathode, copper$^\dagger$ & copper & 9\,\textmu m &  & 0.0627 & 3.16 & 1.66 & 0.85 \\
  drift gas (28.2 mm measured) & Ar/iC$_4$H$_{10}$ 95/5 & 28.2\,mm &  & 0.0236 & 1.83 & 0.96 & 0.49 \\
  \hline
  mesh, steel (calendered) & steel (AISI 316) & 20\,\textmu m & 0.44 & 0.0512 & 2.82 & 1.48 & 0.76 \\
  amplification gap & Ar/iC$_4$H$_{10}$ 95/5 & 150\,\textmu m &  & 0.0001 & 0.09 & 0.05 & 0.03 \\
  resistive strips$^\dagger$ & resistive paste & 10\,\textmu m & 0.20 & 0.0007 & 0.26 & 0.14 & 0.07 \\
  insulator, polyimide + prepreg & glass-epoxy & 50\,\textmu m &  & 0.0258 & 1.93 & 1.01 & 0.52 \\
  pads, copper & copper & 34\,\textmu m & 0.76 & 0.1800 & 5.65 & 2.96 & 1.52 \\
  prepreg & glass-epoxy & 100\,\textmu m &  & 0.0515 & 2.83 & 1.49 & 0.76 \\
  Y strips, copper & copper & 17\,\textmu m & 0.46 & 0.0542 & 2.91 & 1.53 & 0.79 \\
  prepreg & glass-epoxy & 100\,\textmu m &  & 0.0515 & 2.83 & 1.49 & 0.76 \\
  X strips, copper & copper & 17\,\textmu m & 0.46 & 0.0541 & 2.91 & 1.53 & 0.78 \\
  prepreg & glass-epoxy & 100\,\textmu m &  & 0.0515 & 2.83 & 1.49 & 0.76 \\
  via lands, copper & copper & 34\,\textmu m & 0.53 & 0.1256 & 4.63 & 2.43 & 1.25 \\
  prepreg + polyimide & glass-epoxy & 50\,\textmu m &  & 0.0258 & 1.93 & 1.01 & 0.52 \\
  Rohacell$^\dagger$ & Rohacell & 5\,mm &  & 0.0650 & 3.22 & 1.69 & 0.87 \\
  back foil, Mylar$^\dagger$ & Mylar & 25\,\textmu m &  & 0.0088 & 1.06 & 0.55 & 0.29 \\
  \hline
  in front of the drift gas & & & & 0.1712 & 5.49 & 2.88 & 1.48 \\
  \quad without the cathode copper & & & & 0.1085 & 4.27 & 2.24 & 1.15 \\
  including the drift gas & & & & 0.1948 & 5.89 & 3.09 & 1.59 \\
  behind the drift gas & & & & 0.7458 & 12.31 & 6.46 & 3.32 \\
  \hline
  \end{tabular}
\end{table}


\subsection{Production at Saclay}\label{app:process}
\emph{Resistive strips.} The resistive strips are screen printed on the Kapton of the board and cured. Two pastes were
used during the production. The first, from Saral, is soft. It is easy to remove, which matters when a board has to be
printed again, but it did not hold at the step from the copper HV contact down to the board. Strips printed with it
broke there, and the contact had to be repaired with silver glue. Silver is known to migrate under an electric field,
so such a repair may not last. The second paste, from ESL, is hard. It held at the contact without silver glue, and it
can be sanded (800 grit, in both directions) to remove resistive spikes, which cause discharges, and to help the
pillars attach. It is, however, brittle and very hard to remove. All five chambers of this paper have ESL strips. The
pastes and the uniform, strip and grid patterns were compared on small boards
($128\times154$\,mm$^2$) in October and November 2025, with an $^{55}$Fe source, cosmic rays and the SPS beam.
\askStephan{Is there a written result of the small-board tests (uniform/strips/grid, ESL/Saral) we can cite or summarize?}

\emph{Bulk process.} Two layers of Dynamask photoresist (75\,\textmu m each) are laminated on the board, the mesh is laid
on them, and a third layer is laminated over it. The pillars are exposed through a mask with UV light for 20 minutes, so
that the layer next to the Kapton polymerizes fully under the mesh. The board is then developed with a chemical spray from
below, rinsed in de-ionized water, and cured at 145\,\textdegree C, reached in steps of 10\,\textdegree C. Three problems
appeared in this process:
\begin{itemize}
  \item \emph{Waves in the mesh.} At the nominal tension of about 8\,N/cm the mesh formed waves when the board was
  heated for the cure. A tension of about 2\,N/cm avoids them but is hard to reach on the stretching bench, and
  12\,N/cm also avoids them but loads the pillars more. All five chambers were made at 12\,N/cm with the stepped cure.
  Even so, patches of 1--10\,cm across where the mesh is not held by its pillars appeared on several boards, and we do
  not fully understand them.
  \item \emph{Pillars.} Small pillars hide less of the gap but hold less well. A wrinkle in the photoresist during the
  lamination can cost pillars. Before the mesh is laid, the lamination can be removed and redone. After the mesh is
  laid, a wrinkle only costs the cap of a pillar above the mesh, and the pillar under the mesh still holds the gap. The
  development time was kept short to avoid losing pillars. On the last boards this left photoresist under the mesh in
  bands, along the path of the spray, where the efficiency was about 50\,\%.
  \item \emph{Mesh potential.} After the gamma flash of the first n\_TOF test, the mesh was separated from the ground of
  the board, so that a circuit could drive it. This was not part of the design. The mesh was cut on a layer of Kapton
  tape all around the active area and covered with a second layer, and the ground was brought to the mesh through the
  circuit.
\end{itemize}
\askStephan{Pillars as built (\O\,500\,\textmu m on 5\,mm, or \O\,0.6\,mm on 4.68\,mm from the Gerber)? Which of the five
chambers had the developing residue? Is the mesh of chambers A--E floating or grounded at the board?}

\emph{Qualification.} Each board was tested in a class 100\,000 clean room. The aim was 900\,V on the resistive strips
with a current below 1\,\textmu A. A thermal camera located the hot spots, and boards were washed to remove residues of
photoresist. The chambers were then assembled in the clean room and tested on the cosmic bench.

\subsection{The five chambers}\label{app:chambers}
\Cref{tab:production} lists what was seen on each chamber when it was made, next to what we measure. Most boards were
printed and bulked more than once during 2026, and we do not have a record of every iteration.

\begin{table}[tbp]
  \centering\small
  \caption{The five chambers: when the bulk was made, what was seen at production, and what we measure on the bench.}
  \label{tab:production}
  \begin{tabular}{p{0.07\linewidth}p{0.12\linewidth}p{0.38\linewidth}p{0.33\linewidth}}
    \hline
    chamber & bulk made & seen at production & measured \\
    \hline
    A & 15 June & no defect noted & uniform \\
    B & 12 June & pillars too large on parts of the board (the mask shifted during the UV exposure) & uniform \\
    C & about 24 June & lamination redone after wrinkles; a column of waves in the mesh, a spot without pillar caps & discharges over the whole area \\
    D & 22 June & one or two bubbles in the lamination, pillar caps missing there & discharges over the whole area \\
    E & 22 June & a few bubbles in the lamination, pillar caps missing there & bands of broken resistive strips (\cref{app:beam}) \\
    \hline
  \end{tabular}
\end{table}
\askDylan{Chamber D's poorly mounted readout card is an assembly problem, not a bulk one, so it is left out of this
table. Add it? Also: you think C, D and E all have some broken resistive strips; we only show it for E.}
